\documentclass{article}
\usepackage[utf8]{inputenc}
\usepackage[T2A]{fontenc}
\usepackage[russian]{babel}
\usepackage{amsmath}

\begin{document}

\section*{Двойственная норма}

Для нормы $\|\cdot\|$ на $\mathbb{R}^d$ двойственная норма определяется как:
\[
\|g\|_* = \sup_{\|v\|\leq 1}\langle g, v\rangle
\qquad \text{или эквивалентно} \qquad
\|g\|_* = \sup_{v \neq 0} \frac{\langle g, v\rangle}{\|v\|}.
\]

\textbf{Смысл.} Вектор $g$ задаёт линейный функционал $v \mapsto \langle g, v\rangle$.
Двойственная норма $\|g\|_*$ измеряет максимальное значение этого функционала
на единичном шаре по $\|\cdot\|$, то есть насколько сильно $g$ может
<<растянуть>> единичный вектор.

\textbf{Примеры:}
\begin{itemize}
    \item $\ell_2$ самодвойственна: $\|g\|_{2,*} = \|g\|_2$.
    \item $\ell_1$ и $\ell_\infty$ двойственны друг к другу:
          $\|g\|_{1,*} = \|g\|_\infty$ и $\|g\|_{\infty,*} = \|g\|_1$.
    \item $\ell_p$ и $\ell_q$ двойственны при $\tfrac{1}{p}+\tfrac{1}{q}=1$.
\end{itemize}

\textbf{Связь с липшицевостью.} Если $h$ является $L$-липшицевой по $\|\cdot\|$,
то для любого субградиента $g \in \partial h(x_0)$ выполняется $\|g\|_* \leq L$.

\end{document}
